\documentclass[border=12pt,tikz]{standalone}
\usepackage{amsmath,amssymb}
\usepackage{CJKutf8}
\usetikzlibrary{arrows.meta,calc,patterns}

\begin{document}
\begin{CJK}{UTF8}{ipxm}
\begin{tikzpicture}[
    >=Stealth,
    font=\small,
    every node/.style={inner sep=2pt}
]

% --- Parameters ---
\def\R{3.5}
\def\planeW{5.0}

% ============================================================
%  LEFT: R > 0 case (main figure)
% ============================================================

% --- Lower hemisphere (gray, not covered) ---
\draw[thick, gray!40] (0,0) +(-180:{\R}) arc (-180:0:{\R});
\node[gray!50, font=\footnotesize] at (0,-2.2) {（写像の像に含まれない）};

% --- Upper hemisphere (black, covered) ---
\draw[very thick] (0,0) +(0:{\R}) arc (0:180:{\R});

% --- Equator line ---
\draw[thick, gray!60, densely dotted] (-\R,0) -- (\R,0);
\node[gray!60, font=\footnotesize, right=2pt] at (\R,0) {赤道 $Y^{n+1}=0$};

% --- Origin O ---
\filldraw (0,0) circle (2pt);
\node[below=8pt] at (0,0) {$O$（投影中心）};

% --- Contact point P ---
\filldraw (0,\R) circle (2pt);
\node[above left=3pt] at (0,\R) {接点 $P = Re_{n+1}$};

% --- Distance brace R ---
\draw[decorate, decoration={brace, amplitude=5pt, mirror, raise=6pt}]
    (0.05,0.05) -- (0.05,\R-0.05);
\node[right=14pt] at (0,\R/2) {$R > 0$};

% --- Projection center axis ---
\draw[thick, dashed] (0,0) -- (0,\R);
\draw[thick, dashed, -{Stealth[length=6pt]}] (0,\R) -- (0,\R+2.0);
\node[left=4pt] at (0,\R+1.7) {$e_{n+1}$};
\node[right=4pt, font=\footnotesize, blue!60!black] at (0,\R+1.7)
    {（投影中心軸）};

% --- Tangent hyperplane Pi_R ---
\draw[very thick, blue] (-\planeW,\R) -- (\planeW,\R);
\node[blue, below left=2pt] at (-\planeW+0.2,\R) {$\Pi_R$（接超平面）};

% --- Coordinate arrow on Pi_R ---
\draw[-{Stealth[length=5pt]}, blue!60!black, thick]
    (0.6,\R+0.15) -- (4.2,\R+0.15);
\node[blue!60!black, above=1pt] at (2.4,\R+0.15) {座標 $x^1, \ldots, x^n$};

% --- A point on Pi_R ---
\coordinate (Q) at (3.2,\R);
\filldraw[red!80!black] (Q) circle (2.5pt);
\node[above=8pt, red!80!black] at (Q) {$x$};

% --- Compute intersection of ray O->Q with sphere ---
\pgfmathsetmacro{\Qx}{3.2}
\pgfmathsetmacro{\Qy}{\R}
\pgfmathsetmacro{\raylen}{sqrt(\Qx*\Qx + \Qy*\Qy)}
\pgfmathsetmacro{\Sx}{\R*\Qx/\raylen}
\pgfmathsetmacro{\Sy}{\R*\Qy/\raylen}

% --- Projection ray ---
\draw[red!60!black, dashed] (0,0) -- (Q);

% --- Point on sphere ---
\filldraw[red!40!blue] (\Sx,\Sy) circle (2.5pt);
\node[below right=3pt, red!40!blue] at (\Sx,\Sy) {$\Phi(x)$};

% --- Phi arrow ---
\draw[-{Stealth[length=5pt]}, red!40!blue, thick, shorten >=5pt, shorten <=5pt]
    ($(Q)+(0,-0.2)$) to[bend right=25]
    node[left=2pt, pos=0.5, font=\footnotesize]{$\Phi$}
    (\Sx,\Sy);

% --- Annotation: upper hemisphere ---
\node[font=\footnotesize, align=center] at (-3.5,2.0)
    {上半球 $H^+$\\（$Y^{n+1} > 0$）\\$\Phi$ の像};

% --- Background space label ---
\node[gray, font=\normalsize] at (-4.0,-3.5)
    {背景空間 $\mathbb{R}^{n+1}$};

% ============================================================
%  Annotation: R < 0 case (small, below)
% ============================================================

% --- R < 0 tangent plane (dashed, below) ---
\draw[thick, blue!40, dashed] (-3.5,-\R) -- (3.5,-\R);
\node[blue!40, font=\footnotesize, right=2pt] at (3.5,-\R)
    {$R < 0$: $\Pi_R$ は $-e_{n+1}$ 側};

% --- R < 0 contact point ---
\filldraw[gray!60] (0,-\R) circle (2pt);
\node[below left=2pt, gray!60, font=\footnotesize] at (0,-\R) {$P' = Re_{n+1}$（$R<0$）};

% --- Arrow for R < 0 ---
\draw[thick, gray!50, dashed] (0,0) -- (0,-\R);

% --- Annotation ---
\node[gray!70, font=\footnotesize, align=left] at (4.2,-2.5)
    {$R < 0$ のとき\\下半球 $H^-$ を被覆};

\end{tikzpicture}
\end{CJK}
\end{document}
